% mo_hinh_sat_samsp_v2.tex
% Biên dịch: lualatex -interaction=nonstopmode mo_hinh_sat_samsp_v2.tex
%            (chạy 2 lần để cập nhật references)
% KHÔNG dùng pdflatex hoặc xelatex — fontspec yêu cầu lualatex.

\documentclass[11pt,a4paper]{article}

\usepackage[a4paper,margin=25mm]{geometry}
\usepackage{fontspec}
\usepackage{amsmath,amssymb,amsthm,mathtools}
\usepackage{unicode-math}
\usepackage{booktabs,tabularx,longtable,array}
\usepackage{enumitem}
\usepackage{microtype}
\usepackage{hyperref}
\usepackage[nameinlink,noabbrev]{cleveref}

\setmainfont{TeX Gyre Pagella}
\setmathfont{TeX Gyre Pagella Math}

\hypersetup{colorlinks=true,linkcolor=black,citecolor=black,urlcolor=black,
  pdftitle={Q-SAT: Mo hinh SAT compact cho SAMSP}}

\newtheorem{theorem}{Dinh ly}[section]
\newtheorem{lemma}[theorem]{Bo de}
\newtheorem{proposition}[theorem]{Menh de}
\theoremstyle{remark}
\newtheorem{remark}[theorem]{Nhan xet}

\crefname{section}{muc}{cac muc}
\crefname{table}{bang}{cac bang}
\crefname{lemma}{bo de}{cac bo de}
\crefname{theorem}{dinh ly}{cac dinh ly}

\newcommand{\AMO}{\operatorname{AMO}}
\newcommand{\ES}{\operatorname{ES}}
\newcommand{\LS}{\operatorname{LS}}
\newcommand{\EC}{\operatorname{EC}}
\newcommand{\LC}{\operatorname{LC}}
\newcommand{\STW}{\operatorname{STW}}
\newcommand{\RTW}{\operatorname{RTW}}
\newcolumntype{Y}{>{\raggedright\arraybackslash}X}
\newcolumntype{Z}{>{\centering\arraybackslash}X}

\setlist[itemize]{leftmargin=1.4em,itemsep=0.2em,topsep=0.3em}
\setlist[enumerate]{leftmargin=1.6em,itemsep=0.2em,topsep=0.3em}

\begin{document}

\title{\textbf{Q-SAT: Mô hình SAT compact cho\\
    Sample Analysis Machine Scheduling Problem}}
\author{Tuyên Văn Kiều}
\date{Tháng 9 năm 2026 \quad\emph{(bản kỹ thuật nội bộ)}}
\maketitle

\begin{abstract}
Bài toán SAMSP kết hợp generalized precedence (kể cả maximum lag), tài
nguyên tái tạo unary và vùng lưu trữ có sức chứa hữu hạn.
Chúng tôi đề xuất Q-SAT: mô hình SAT time-indexed dùng duy nhất
\emph{order variables} $q_{i,t} \equiv [S_i \le t]$ để biểu diễn đồng thời
start time, precedence và storage, loại bỏ start literals và register variables
của baseline hiện tại.
Q-SAT được chứng minh tương đương nghiệm với các mô hình CP, SMT và MILP
của bài báo gốc.
Ngoài ra, bài viết làm rõ điều kiện áp dụng PB(AMO) cho generalized lags
và phân biệt rõ ladder encoding với SCAMO.
\end{abstract}

\tableofcontents
\clearpage

%=============================================================================
\section{Định nghĩa SAMSP}
\label{sec:samsp}
%=============================================================================

\subsection{Dữ liệu bài toán}

Một instance SAMSP được cho bởi bộ
$\mathcal{T} = (\mathcal{V}, d, \mathcal{E}, \mathcal{R}, B, b, \mathcal{C}, P, c)$,
trong đó:

\begin{itemize}
  \item $\mathcal{V} = \{A_0, A_1, \ldots, A_n, A_{n+1}\}$: tập hoạt động;
    $A_0, A_{n+1}$ là dummy start/end với $d_0 = d_{n+1} = 0$.
  \item $d_i \in \mathbb{N}$: thời lượng của $A_i$.
  \item $(A_i, A_j, \ell_{ij}) \in \mathcal{E}$: generalized precedence
    $S_j - S_i \ge \ell_{ij}$, với $\ell_{ij}$ có thể âm.
    Maximum lag $S_j - S_i \le L$ được mã hóa bằng cạnh ngược
    $S_i - S_j \ge -L$.
  \item $\mathcal{R}$: tập tài nguyên tái tạo; $B_r$: sức chứa;
    $b_{ir} \ge 0$: nhu cầu của $A_i$ trong suốt thời gian chạy.
  \item $\mathcal{C}$: tập vùng lưu trữ; $P_s$: sức chứa.
  \item $c_{is} \in \mathbb{Z}$: thay đổi mức storage ---
    $c_{is} > 0$ chiếm chỗ ngay khi $A_i$ bắt đầu;
    $c_{is} < 0$ trả $|c_{is}|$ chỗ ngay sau khi $A_i$ hoàn tất.
\end{itemize}

Mức chiếm storage $s$ tại biên thời gian $t$:
\begin{equation}
  O_s(t) =
  \sum_{i:\,c_{is}>0} c_{is}\,[S_i \le t]
  + \sum_{i:\,c_{is}<0} c_{is}\,[S_i + d_i \le t].
  \label{eq:storage-level}
\end{equation}

\subsection{Bài toán quyết định}

Với cận makespan $T$, tìm vector start time $S$ thỏa:
\begin{align}
  S_0 &= 0, & S_{n+1} &\le T, \label{eq:dummy}\\
  S_j - S_i &\ge \ell_{ij} && \forall (i,j,\ell_{ij}) \in \mathcal{E},
    \label{eq:prec}\\
  \sum_{i:\,t \in [S_i,\,S_i+d_i-1]} b_{ir} &\le B_r
    && \forall r,\; t, \label{eq:renewable}\\
  0 \le O_s(t) &\le P_s && \forall s,\; t. \label{eq:storage}
\end{align}

\begin{remark}
Khoảng chiếm storage không trùng với duration của một activity:
hai đầu khoảng do hai activities khác nhau quyết định, với độ dài phụ thuộc
vào nhiều lags. Gộp khoảng chiếm thành một activity dài là sai về semantics.
\end{remark}

\subsection{Ví dụ một cuvette}
\label{sec:example}

Xét sáu hoạt động (mỗi renewable resource và storage có capacity 1):

\begin{table}[htbp]
\centering
\caption{Dữ liệu ví dụ một cuvette. $+C/-C$, $+H/-H$, $+O/-O$: chiếm/trả
  Cold, Hot-II, Observation.}
\label{tab:example}
\begin{tabularx}{\textwidth}{@{}cYcp{0.20\textwidth}p{0.20\textwidth}@{}}
\toprule
\textbf{Act.} & \textbf{Thao tác} & $\boldsymbol{d_i}$ &
\textbf{Renewable} & \textbf{Storage event}\\
\midrule
$A_1$ & Đưa cuvette vào Cold   & 2 & Shuttle     & $+C$\\
$A_2$ & Hút mẫu vào cuvette    & 3 & Sample Arm  & ---\\
$A_3$ & Chuyển Cold $\to$ Hot-II   & 2 & Shuttle     & $-C,\;+H$\\
$A_4$ & Thêm thuốc thử         & 2 & Reagent Arm & ---\\
$A_5$ & Chuyển Hot-II $\to$ Obs.   & 2 & Shuttle     & $-H,\;+O$\\
$A_6$ & Chuyển Obs. $\to$ waste    & 2 & Shuttle     & $-O$\\
\bottomrule
\end{tabularx}
\end{table}

Minimum lags: $S_2-S_1\ge2$, $S_3-S_2\ge3$, $S_4-S_3\ge2$,
$S_5-S_4\ge2$, $S_6-S_5\ge6$.
Maximum lag: $S_6-S_5\le8$, mã hóa bằng cạnh ngược $S_5-S_6\ge-8$.
Lịch khả thi: $(S_1,\ldots,S_6,S_7) = (0,2,5,7,9,15,17)$, makespan $= 17$.

%=============================================================================
\section{Các mô hình gốc}
\label{sec:original}
%=============================================================================

Bofill et al.~\cite{bofill2022samsp} xây dựng ba mô hình exact: CP, SMT và MILP.

\begin{table}[htbp]
\centering
\caption{So sánh ba mô hình của bài báo gốc.}
\label{tab:original}
\begin{tabularx}{\textwidth}{@{}p{0.10\textwidth}p{0.25\textwidth}Y Y@{}}
\toprule
\textbf{Mô hình} & \textbf{Biến chính} & \textbf{Capacity constraints} &
\textbf{Đặc điểm}\\
\midrule
CP   & Interval variable $\operatorname{task}_i$ &
  Pulse, StepAtStart, StepAtEnd &
  Không cần horizon để tạo biến; propagation toàn cục.\\
SMT  & $S_i \in \mathbb{Z}$; $X_{i,t}, Y_{i,t} \in \{0,1\}$ &
  IDL cho precedence; PB/AMO cho renewable; PB cho storage &
  Time-indexed ở $X,Y$; start time vẫn là biến nguyên.\\
MILP & One-hot $\xi_{i,t} \in \{0,1\}$ &
  Tổng chập start literals (renewable); tổng tích lũy (storage) &
  Hoàn toàn time-indexed; LP relaxation cung cấp LB.\\
\bottomrule
\end{tabularx}
\end{table}

\paragraph{Ký hiệu dùng lại.}
$\STW(i) = [\ES_i, \LS_i]$: start time window;
$\RTW(i) = [\ES_i, \LC_i - 1]$: run time window.
Mô hình MILP time-indexed (dùng làm baseline):
\begin{alignat}{2}
  \min\quad & \textstyle\sum_t t\,\xi_{n+1,t}, \tag{ILP-OBJ}\\
  & \textstyle\sum_{t \in \STW(i)} \xi_{i,t} = 1 && \forall i,
    \tag{ILP-EO}\label{ILP-EO}\\
  & \textstyle\sum_t t\,\xi_{j,t} - \sum_t t\,\xi_{i,t} \ge \ell^\star_{ij}
    && \forall (i,j) \in \mathcal{E}^\star, \tag{ILP-P}\\
  & \textstyle\sum_i b_{ir} R_{i,t}(\xi) \le B_r
    && \forall r,\;t, \tag{ILP-R}\\
  & \textstyle 0 \le \sum_{i:\,c_{is}>0} c_{is} Q_{i,t}(\xi)
    + \sum_{i:\,c_{is}<0} c_{is} Q_{i,t-d_i}(\xi) \le P_s
    && \forall s,\;t, \tag{ILP-C}
\end{alignat}
với $Q_{i,t} = \sum_{\tau \le t, \tau \in \STW(i)} \xi_{i,\tau}$ và
$R_{i,t} = \sum_{\tau=\max(\ES_i,t-d_i+1)}^{\min(t,\LS_i)} \xi_{i,\tau}$.

%=============================================================================
\section{Preprocessing}
\label{sec:preprocessing}
%=============================================================================

\subsection{Closure max-plus}

$\ell^\star_{ij}$ là lag kéo theo mạnh nhất từ tất cả đường đi $i \leadsto j$:
\[
  \ell^\star_{ij} = \max_{\pi:\,i\leadsto j} \sum_{(a,b)\in\pi} \ell_{ab}.
\]
Tính bằng Floyd--Warshall với phép cập nhật
$L_{ij} \leftarrow \max\{L_{ij},\, L_{ik} + L_{kj}\}$ (\emph{max-plus},
không phải shortest path).
Nếu $L_{ii} > 0$ sau closure, instance UNSAT (positive cycle).

\begin{remark}
\texttt{networkx.floyd\_warshall} tính shortest path.
Để dùng đúng: đặt $w_{ij} = -\ell_{ij}$ trước khi gọi, rồi đổi dấu kết quả;
hoặc cài max-plus trực tiếp.
\end{remark}

\subsection{Time windows}

Cận dưới: $LB_0 = \ell^\star_{0,n+1}$.
Cận trên ban đầu: $UB_0 = \sum_i \gamma_i$ với $\gamma_i = \max(d_i, \max_{(i,j)}\ell_{ij})$.
$UB_0$ thường lỏng; nên thay bằng incumbent từ greedy schedule.

Với $\overline{U} = \min(UB_0, UB_\mathrm{inc})$:
\begin{equation}
  \ES_i = \ell^\star_{0i},\quad \LS_i = \overline{U} - \ell^\star_{i,n+1},\quad
  \EC_i = \ES_i + d_i,\quad \LC_i = \LS_i + d_i.
  \label{eq:windows}
\end{equation}
Nếu $\ES_i > \LS_i$, instance UNSAT không cần gọi solver.

\subsection{Symmetry breaking}
\label{sec:sb}

Với activities exchangeable (cùng test type, mọi dữ liệu đồng nhất):

\begin{description}[style=nextline]
  \item[SB-R (renewable).] Với mỗi cặp copy liên tiếp $(m, m+1)$:
  \[
    (A_{k,m,a},\; A_{k,m+1,a},\; o_a),\quad
    o_a = \begin{cases} d_a & \text{nếu } \sum_r b_{ar} > 0,\\
      0 & \text{ngược lại.}\end{cases}
  \]

  \item[SB-C (storage).] Đặt $h = \lfloor P_s / c_{as} \rfloor$ và $pt$ là
  thời gian tối thiểu từ start của $a$ đến completion của activity trả một vị
  trí $s$ (closure cục bộ test type $k$):
  \[
    (A_{k,m,a},\; A_{k,m+h,a},\; pt),\quad m = 1, \ldots, N_k - h.
  \]
\end{description}

Symmetry edges phải được thêm \emph{trước} closure toàn cục để ảnh hưởng
của chúng lan qua toàn graph.

\paragraph{Pipeline.}
\begin{enumerate}
  \item Closure cục bộ (từng test type) $\to$ sinh SB-R, SB-C.
  \item Closure toàn cục trên $\mathcal{E}_g$ đã bổ sung; phát hiện positive cycle.
  \item Tính $LB_0$; chạy greedy để lấy $\overline{U}$; tính windows \eqref{eq:windows}.
  \item (Tùy chọn) Energetic strengthening: thêm cạnh, lặp closure đến fixed point.
  \item Thay literal ngoài cửa sổ bằng hằng; chọn chính sách materialization.
\end{enumerate}

%=============================================================================
\section{Mô hình Q-SAT}
\label{sec:qsat}
%=============================================================================

\subsection{Order variables}

\begin{equation}
  q_{i,t} = 1 \iff S_i \le t, \qquad t \in [\ES_i, \LS_i - 1].
  \label{eq:order-var}
\end{equation}
Biên cố định (hằng, không phải biến SAT):
\begin{equation}
  q_{i,t} = \begin{cases} 0 & t < \ES_i,\\ 1 & t \ge \LS_i. \end{cases}
  \label{eq:boundary}
\end{equation}

\begin{lemma}[Ánh xạ 1-1]
\label{lem:transition}
Mọi gán thỏa \eqref{eq:boundary} và các mệnh đề đơn điệu
\begin{equation}
  \neg q_{i,t-1} \lor q_{i,t}, \qquad t = \ES_i + 1, \ldots, \LS_i - 1,
  \label{eq:monotone}
\end{equation}
xác định duy nhất $S_i = \min\{t : q_{i,t} = 1\}$, và ngược lại.
\end{lemma}

\begin{proof}
Chuỗi đơn điệu bị chặn bởi 0 ở trái và 1 ở phải, nên có đúng một chuyển
mức $0 \to 1$. Vị trí đó chính là $S_i$.
\end{proof}

Start literal dẫn xuất (không cần tạo biến SAT riêng):
\[
  \sigma_{i,t} \equiv q_{i,t} \land \neg q_{i,t-1} \iff S_i = t.
\]

\subsection{Precedence}

Với cạnh $(i, j, \ell) \in \mathcal{E}^\star$:
\begin{equation}
  \neg q_{j,t} \lor q_{i,t-\ell}, \qquad t \in [\ES_j, \LS_j].
  \label{eq:prec-clause}
\end{equation}

\begin{lemma}
\label{lem:prec}
Các mệnh đề \eqref{eq:prec-clause} tương đương $S_j - S_i \ge \ell$,
với $\ell$ tùy ý (dương hoặc âm).
\end{lemma}

\begin{proof}
Tại $t = S_j$: $q_{j,t}=1$ buộc $q_{i,S_j-\ell}=1$, tức $S_i \le S_j - \ell$.
Chiều ngược: nếu $S_i \le S_j - \ell$ và $q_{j,t}=1$ thì $t \ge S_j$,
suy ra $t-\ell \ge S_i$ và $q_{i,t-\ell}=1$.
\end{proof}

\subsection{Run variables và channeling}

\begin{equation}
  u_{i,t} = 1 \iff t \in [S_i,\; S_i + d_i - 1], \qquad t \in \RTW(i).
  \label{eq:run-var}
\end{equation}

Channeling ($u_{i,t} \iff q_{i,t} \land \neg q_{i,t-d_i}$) qua ba clause:
\begin{align}
  \neg u_{i,t} \lor q_{i,t}, \tag{17a}\label{eq:ch-a}\\
  \neg u_{i,t} \lor \neg q_{i,t-d_i}, \tag{17b}\label{eq:ch-b}\\
  \neg q_{i,t} \lor q_{i,t-d_i} \lor u_{i,t}. \tag{17c}\label{eq:ch-c}
\end{align}

\subsection{Tài nguyên tái tạo}

\begin{equation}
  \sum_{i:\,t \in \RTW(i)} b_{ir}\, u_{i,t} \le B_r,
  \qquad \forall r,\; t.
  \label{eq:renewable-pb}
\end{equation}
Mã hóa sang CNF bằng BDD~\cite{abio2012bdd}, AMO (khi $B_r=1$),
hoặc PB(AMO) (\cref{sec:pbamo}).

\subsection{Storage}

Đặt $I_s^+ = \{i : c_{is} > 0\}$, $I_s^- = \{i : c_{is} < 0\}$,
$w_{is} = -c_{is} > 0$.
Từ \eqref{eq:storage-level} và \eqref{eq:order-var}:
\begin{equation}
  O_s(t) = \sum_{i \in I_s^+} c_{is}\, q_{i,t}
          - \sum_{i \in I_s^-} w_{is}\, q_{i,t-d_i}.
  \label{eq:storage-q}
\end{equation}
Quy ước biên \eqref{eq:boundary} hấp thụ tự động phần level đã cố định
của mô hình SMT gốc.

Cận trên $O_s(t) \le P_s$ (sau khi chuyển về trọng số dương):
\begin{equation}
  \sum_{i \in I_s^+} c_{is}\, q_{i,t}
  + \sum_{i \in I_s^-} w_{is}\, \neg q_{i,t-d_i}
  \le P_s + \textstyle\sum_{i \in I_s^-} w_{is}.
  \label{eq:storage-ub}
\end{equation}

Cận dưới $O_s(t) \ge 0$:
\begin{equation}
  \sum_{i \in I_s^-} w_{is}\, q_{i,t-d_i}
  + \sum_{i \in I_s^+} c_{is}\, \neg q_{i,t}
  \le \textstyle\sum_{i \in I_s^+} c_{is}.
  \label{eq:storage-lb}
\end{equation}

\paragraph{Ví dụ (hai cuvettes, $P_s = 1$).}
Xét hai cuvettes $p_1, p_2$ chiếm Cold ($c=+1$) và $v_1, v_2$ trả Cold
($c=-1$, $d_v=2$). Tại biên $t$:
\[
  O_C(t) = q_{p_1,t} + q_{p_2,t} - q_{v_1,t-2} - q_{v_2,t-2}.
\]
Cận trên \eqref{eq:storage-ub}:
\[
  q_{p_1,t} + q_{p_2,t} + \neg q_{v_1,t-2} + \neg q_{v_2,t-2} \le 1 + 2 = 3.
\]
Vi phạm xảy ra khi cả hai cuvettes đã vào Cold mà chưa cuvette nào được trả:
$q_{p_1,t}=q_{p_2,t}=1$, $q_{v_1,t-2}=q_{v_2,t-2}=0$, cho vế trái $= 4 > 3$.

\subsection{Makespan và giải incremental}

Assumption $q_{n+1,T}$ truy vấn $S_{n+1} \le T$.
Chiến lược tối ưu hóa:
\begin{enumerate}
  \item Binary search trên $[LB_0, \overline{U}]$: $O(\log(\overline{U}-LB_0))$ SAT calls.
  \item Top-down descent từ $T^\star$: learned clauses từ SAT$(T^\star)$
    được tái sử dụng cho SAT$(T^\star-1)$.
\end{enumerate}
Giải mã lịch: $S_i = \min\{t : q_{i,t} = 1\}$ (\cref{lem:transition}).

\subsection{Công thức Q-SAT hoàn chỉnh}

\begin{equation}
\mathcal{F}_T =
  \bigwedge_{i,t} \eqref{eq:monotone}
  \;\wedge\; \bigwedge_{(i,j,\ell),\,t} \eqref{eq:prec-clause}
  \;\wedge\; \bigwedge_{i,t} \eqref{eq:ch-a}\text{--}\eqref{eq:ch-c}
  \;\wedge\; \bigwedge_{r,t} \operatorname{Enc}\eqref{eq:renewable-pb}
  \;\wedge\; \bigwedge_{s,t} \operatorname{Enc}\!\left[\eqref{eq:storage-ub}\wedge\eqref{eq:storage-lb}\right]
  \;\wedge\; q_{n+1,T}.
\tag{Q-SAT}
\label{eq:qsat}
\end{equation}

\begin{theorem}[Soundness và completeness]
\label{thm:correct}
Giả sử time windows chứa mọi lịch có makespan $\le T$ và
$\operatorname{Enc}$ là đúng. Khi đó $\mathcal{F}_T$ satisfiable khi và chỉ
khi SAMSP có lịch khả thi với makespan $\le T$.
\end{theorem}

\begin{proof}[Phác thảo]
($\Leftarrow$) Đặt $q_{i,t} = [S_i \le t]$, $u_{i,t} = [S_i \le t < S_i+d_i]$.
\Cref{lem:transition,lem:prec} và \eqref{eq:ch-a}--\eqref{eq:ch-c} đảm bảo core clauses đúng.
\eqref{eq:storage-q} cho $O_s(t)$ đúng, nên \eqref{eq:storage-ub}--\eqref{eq:storage-lb} đúng.

($\Rightarrow$) \Cref{lem:transition} giải mã duy nhất $S_i$;
\cref{lem:prec} cho precedence; \eqref{eq:ch-a}--\eqref{eq:ch-c} cho $u_{i,t}$
đúng trong $[S_i, S_i+d_i-1]$; PB encoder cho \eqref{eq:renewable}--\eqref{eq:storage}.
\end{proof}

%=============================================================================
\section{Ánh xạ với SMT, MILP và CP}
\label{sec:map}
%=============================================================================

\begin{table}[htbp]
\centering
\caption{Ánh xạ các biến của bốn mô hình.}
\label{tab:map}
\begin{tabularx}{\textwidth}{@{}p{0.30\textwidth}p{0.22\textwidth}Y@{}}
\toprule
\textbf{Đối tượng trong mô hình gốc} & \textbf{Q-SAT} & \textbf{Ngữ nghĩa}\\
\midrule
$S_i$ (SMT); $\sum_t t\xi_{i,t}$ (MILP) & $\min\{t : q_{i,t}=1\}$ & Start time.\\
$Y_{i,t}$ (SMT); $Q_{i,t}$ (MILP); StepAtStart (CP) & $q_{i,t}$ & Đã bắt đầu $\le t$.\\
$X_{i,t}$ (SMT); $R_{i,t}$ (MILP); Pulse (CP) & $u_{i,t}$ & Đang chạy tại $t$.\\
$\xi_{i,t}$ (MILP) & $q_{i,t} \land \neg q_{i,t-1}$ & Bắt đầu đúng tại $t$.\\
StepAtEnd (CP); $Y_{i,t-d_i}$ (SMT) & $q_{i,t-d_i}$ & Đã hoàn tất $\le t$.\\
\bottomrule
\end{tabularx}
\end{table}

\paragraph{Quan hệ chính.}
$Y_{i,t} \equiv q_{i,t}$ và $X_{i,t} \equiv u_{i,t}$.
Phần level cố định $PS_{s,t}$ của SMT tương đương phần hằng trong
\eqref{eq:storage-q} do quy ước biên \eqref{eq:boundary}.
$Q_{i,t}(\xi) \equiv q_{i,t}$; $R_{i,t}(\xi) \equiv u_{i,t}$ dưới exactly-one.

\begin{theorem}[Tương đương nghiệm]
\label{thm:equiv}
Bốn mô hình CP, SMT, MILP và Q-SAT biểu diễn cùng tập lịch khả thi
$\mathcal{F}_\mathrm{SAMSP}$.
\end{theorem}

\begin{proof}
Dùng các đẳng thức trong \cref{tab:map}: từ mỗi $S \in \mathcal{F}_\mathrm{SAMSP}$,
gán thoả đồng thời bốn mô hình; chiều ngược lại, mỗi model xác định $S$
thỏa \eqref{eq:dummy}--\eqref{eq:storage}.
\end{proof}

\begin{remark}
Tương đương về tập nghiệm không kéo theo tương đương về suy diễn.
CP, SMT, MILP và SAT duyệt không gian tìm kiếm rất khác nhau.
\end{remark}

%=============================================================================
\section{PB(AMO) từ extended precedence}
\label{sec:pbamo}
%=============================================================================

\subsection{Non-overlap graph}

\begin{equation}
  D = (\mathcal{A}, F),\qquad (i,j) \in F \iff \ell^\star_{ij} \ge d_i.
  \label{eq:nooverlap}
\end{equation}
$(i,j) \in F$ chứng minh $A_i$ hoàn tất trước khi $A_j$ bắt đầu, nên
$u_{i,t}$ và $u_{j,t}$ không thể đồng thời bằng 1.

Path cover đỉnh-rời nhau $\mathcal{P}_{r,t}$ của $D[V_{r,t}]$
với $V_{r,t} = \{i : b_{ir} > 0,\, t \in \RTW(i)\}$ cho:
\begin{equation}
  \sum_{i \in P_h} u_{i,t} \le 1, \qquad h = 1, \ldots, |\mathcal{P}_{r,t}|.
  \label{eq:amo-path}
\end{equation}
Kết hợp \eqref{eq:amo-path} với \eqref{eq:renewable-pb} tạo PB(AMO)~\cite{bofill2021pbamo}.
Minimum path cover tính qua maximum bipartite matching.

\begin{remark}
Path cover trên $D$ không suy ra AMO trên storage literals $q_{i,t}$,
vì nhiều $q_{i,t}$ trên cùng đường đi có thể đồng thời bằng 1
(chúng biểu diễn ``đã bắt đầu'', không phải ``đang chạy'').
\end{remark}

\subsection{Selective materialization}

Cấu hình selective chỉ đưa vào CNF các cạnh hữu ích cho cả precedence
clauses lẫn non-overlap graph:
\[
  \mathcal{E}_\mathrm{sel} =
  \mathcal{E}_\mathrm{raw}
  \cup \mathcal{E}_\mathrm{SB}
  \cup \{(i,j,\ell^\star_{ij}) : \ell^\star_{ij} \ge d_i\},
\]
trong khi vẫn dùng toàn bộ closure để tính windows.

%=============================================================================
\section{Ladder, Staircase và SCAMO}
\label{sec:staircase}
%=============================================================================

Chuỗi $q_{i,\cdot}$ là \emph{ladder} đơn điệu, mã hóa bằng binary clause
\eqref{eq:monotone}. Đây \emph{không phải} SCAMO.
SCAMO($\Omega$, $w$) đòi mọi cửa sổ kích thước $w$ là AMO; chuỗi $q_i$
không có cấu trúc cửa sổ đó.

\begin{table}[htbp]
\centering
\caption{Encoding phù hợp cho từng cấu trúc trong SAMSP.}
\label{tab:scamo}
\begin{tabularx}{\textwidth}{@{}p{0.30\textwidth}p{0.14\textwidth}Y@{}}
\toprule
\textbf{Cấu trúc} & \textbf{SCAMO?} & \textbf{Encoding phù hợp}\\
\midrule
Chuỗi $q_{i,t}$ đơn điệu & Không (ladder) & Clause \eqref{eq:monotone}.\\
No-overlap $\{u_{i,t}\}$ của copies cùng resource, trượt theo $t$
  & Có & SCL~\cite{truong2025scl} hoặc order clauses \eqref{eq:prec-clause}.\\
Resource $B_r = 1$, tập $V_{r,t}$ thay đổi theo $t$
  & AMO tại mỗi $t$ & AMO riêng hoặc path-cover PB(AMO).\\
Resource $B_r > 1$ & Không & PB/cardinality encoding.\\
Storage \eqref{eq:storage-ub}--\eqref{eq:storage-lb} & Không (chung)
  & PB encoding.\\
\bottomrule
\end{tabularx}
\end{table}

%=============================================================================
\section{Phân tích kích thước}
\label{sec:size}
%=============================================================================

$W_i = \LS_i - \ES_i + 1$; $U_i = W_i + d_i - 1$.
Không tính biến phụ của PB encoders:

\begin{table}[htbp]
\centering
\caption{Kích thước phần temporal (trước PB encoding).}
\label{tab:size}
\begin{tabularx}{\textwidth}{@{}Y Z Z Y@{}}
\toprule
\textbf{Thành phần} & \textbf{Baseline} & \textbf{Q-SAT} & \textbf{Nhận xét}\\
\midrule
Biến temporal/activity & $2W_i + U_i$ & $(W_i-1) + U_i$ & Tiết kiệm $W_i+1$ biến.\\
Clause đơn điệu & $\approx 4W_i - 1$ & $\max(0, W_i-2)$ & Sau khi rút gọn hằng.\\
Clause start-to-run & $W_i d_i$ & 0 & Thay bởi 3 channeling clauses.\\
Clause run channel & $\approx W_i - 2$ & $\le 3U_i$ & Semantics iff.\\
Biến storage "started" & --- / $Y_{i,t}$ (SMT) & 0 thêm & Tái sử dụng $q$.\\
\bottomrule
\end{tabularx}
\end{table}

Biến/clause PB thường lớn hơn phần temporal và phụ thuộc encoder.
Báo cáo tách riêng $(\#V_\mathrm{core}, \#C_\mathrm{core}, \#V_\mathrm{PB}, \#C_\mathrm{PB})$.

%=============================================================================
\section{Thiết kế thực nghiệm}
\label{sec:experiment}
%=============================================================================

\paragraph{Dữ liệu.}
Ba tập benchmark:
(i) SAMSP-Original: 48 instance của~\cite{bofill2022samsp} (cần liên hệ tác giả);
(ii) RCPSP-Regression: PSPLIB $j30$/$j60$ để kiểm tra lõi Q-SAT;
(iii) SAMSP-Synthetic: sinh có seed, thay đổi số copies, storage capacity,
lag density, tỷ lệ maximum lag; bao gồm cả UNSAT instances.

\paragraph{Ablation.}

\begin{table}[htbp]
\centering
\caption{Ma trận ablation.}
\label{tab:ablation}
\begin{tabularx}{\textwidth}{@{}p{0.16\textwidth}Y Y@{}}
\toprule
\textbf{Nhóm} & \textbf{Cấu hình so sánh} & \textbf{Câu hỏi}\\
\midrule
Temporal core & Baseline / Q-SAT & Giảm biến có giảm conflicts?\\
Precedence & Raw / +SB / bounds-only / selective / full & Nên materialize closure như thế nào?\\
Resource PB & BDD / AMO / PB(AMO) & Encoder nào thắng theo capacity và horizon?\\
Storage & Cả hai cận / chỉ cận trên & Khi nào bỏ cận dưới an toàn?\\
Optimization & Rebuild / binary / binary+top-down & Learned clause reuse bù được windows lỏng?\\
\bottomrule
\end{tabularx}
\end{table}

\paragraph{Chỉ số báo cáo.}
Theo nhóm: $\#V_\mathrm{core}$, $\#C_\mathrm{core}$, $\#V_\mathrm{PB}$,
$\#C_\mathrm{PB}$, thời gian preprocessing/encoding/solving, số SAT calls,
decisions, conflicts, propagations, makespan, LB certified.
Dùng geometric mean và performance profile.

\paragraph{Giả thuyết chính.}
\begin{enumerate}
  \item Q-SAT giảm rõ rệt core variables/clauses, nhất là khi $W_i d_i$ lớn.
  \item Full closure tốt khi windows co mạnh; thua selective khi $|\mathcal{E}^\star|$ lớn.
  \item PB(AMO) có lợi khi path cover nhỏ; không có lợi khi lag cho phép nhiều overlap.
  \item Bỏ cận dưới storage chỉ an toàn sau khi validator xác nhận dữ liệu well-formed.
\end{enumerate}

%=============================================================================
\section{Hướng cải tiến tiếp theo}
\label{sec:future}
%=============================================================================

\begin{table}[htbp]
\centering
\caption{Các hướng nghiên cứu sau Q-SAT.}
\label{tab:future}
\begin{tabularx}{\textwidth}{@{}p{0.21\textwidth}Y Y@{}}
\toprule
\textbf{Hướng} & \textbf{Lợi thế} & \textbf{Rủi ro}\\
\midrule
Energetic Reasoning LB
  & $LB = \max(LB_\mathrm{CPM}, LB_\mathrm{ER}, LB_\mathrm{storage})$,
    thu hẹp binary search
  & Cần xử lý storage với hệ số âm.\\
Primal heuristic greedy
  & $\overline{U}$ nhỏ $\to$ windows hẹp $\to$ formula nhỏ
  & Greedy bị chặn bởi maximum lag.\\
MaxSAT trực tiếp (RC2/EvalMaxSAT)
  & Không cần binary search
  & Hiện tại yếu hơn incremental SAT trên bài toán lớn.\\
Probing / failed literal
  & Thu hẹp windows sau closure; miễn phí với PySAT
  & Tốn kém nếu không có termination condition.\\
Paired occupancy intervals
  & Storage trở thành cumulative/cardinality đơn giản
  & Pairing phải được suy ra chắc chắn từ dữ liệu.\\
\bottomrule
\end{tabularx}
\end{table}

%=============================================================================
\section{Sinh dataset SAMSP}
\label{sec:dataset}
%=============================================================================

\subsection{SAMSP cần gì thêm so với RCPSP}

Một instance RCPSP gồm $n$ activities, $K$ tài nguyên, duration $d$,
demand $b$, capacity $B$ và DAG precedence với lag end-start.
SAMSP mở rộng ba điểm:

\begin{enumerate}
  \item \textbf{Generalized lags:} $\ell_{ij} \in \mathbb{Z}$ (âm hoặc dương);
    maximum lag $S_j - S_i \le L$ mã hóa bằng cạnh ngược
    $(j, i, -L)$.
  \item \textbf{Storage zones:} $S$ storages với sức chứa $P_s$;
    storage event $c_{is} \in \mathbb{Z}$ cho từng cặp (activity, storage).
  \item \textbf{Storage balance:} mỗi ``cuvette'' phải thỏa
    $\sum_i c_{is} = 0$ (chiếm bao nhiêu trả bấy nhiêu) và
    $0 \le O_s(t) \le P_s$ tại mọi $t$.
\end{enumerate}

\subsection{Bốn hướng sinh instance}

\paragraph{Hướng 1 — Template cố định.}
Mã hóa cứng 4 test type từ bài báo gốc; mỗi instance là tổ hợp số copies
$\langle N_1, N_2, N_3, N_4 \rangle$. Phù hợp để tái lập kết quả bài báo
gốc, nhưng không đủ đa dạng cho ablation.

\paragraph{Hướng 2 — Parameterized template (đề xuất chính).}
Mỗi test type được sinh từ tập tham số:
\begin{itemize}[noitemsep]
  \item $n_\mathrm{acts}$: số activities/type; $[d_{\min}, d_{\max}]$: phân phối duration.
  \item $\lambda_\mathrm{lag}$: mật độ lag (xác suất thêm cạnh shortcut).
  \item $\rho_\mathrm{max}$: tỷ lệ min lags cũng có max lag.
  \item $\delta_\mathrm{slack}$: $L = \ell_{\min} \cdot (1 + \delta)$,
    $\delta \in [0.1,\; 3.0]$.
  \item $\sigma_\mathrm{storage}$: tỷ lệ activities có storage event.
  \item $[P_{\min}, P_{\max}]$: phân phối storage capacity.
\end{itemize}
Instance được ghép từ nhiều copies của một hoặc nhiều test types.

\paragraph{Hướng 3 — ProGen-inspired.}
Kế thừa 4 tham số của ProGen (NC, RS, RF, OS) và thêm 4 tham số SAMSP:
\begin{itemize}[noitemsep]
  \item \textbf{SF} (storage factor): tỷ lệ activities có storage event.
  \item \textbf{SS} (storage strength): $(P_s - 1) / (P_{s,\max} - 1)$
    — nghịch với độ chặt storage.
  \item \textbf{LF} (lag factor): tỷ lệ cạnh có max lag.
  \item \textbf{LS} (lag slack): $L / \ell_{\min} - 1$.
\end{itemize}
Framework này cho phép so sánh với RCPSP literature theo cùng không gian tham số.

\paragraph{Hướng 4 — Hybrid library.}
Xây thư viện 10--20 test type templates (4 từ paper gốc + phần còn lại
sinh bằng Hướng 2); mỗi instance chọn ngẫu nhiên subset và số copies.
Đây là hướng phù hợp nhất cho publication vì vừa đa dạng vừa có cấu trúc.

\subsection{Vấn đề trọng tâm: storage balance}

Sinh $c_{is}$ ngẫu nhiên độc lập sẽ vi phạm global balance hoặc tạo level
âm. Giải pháp: sinh storage events \emph{theo cặp}.

\begin{verbatim}
for s in storages:
    for _ in range(n_cuvettes):
        # Chọn put_act: activity đặt cuvette vào storage s
        put = choice([a for a in acts if can_put(a, s)])
        # Chọn vacate_act: phải đến sau put theo topological order
        vac = choice([a for a in acts
                      if topo_rank[a] > topo_rank[put]
                      and can_vacate(a, s)])
        c[put][s] += 1    # chiếm tại start của put
        c[vac][s] -= 1    # trả sau completion của vac
\end{verbatim}

Sau khi sinh, chạy một greedy schedule và kiểm tra
$O_s(t) \in [0, P_s]$ tại mọi $t$ (\cref{eq:storage}).
Nếu vi phạm, hoán đổi một số cặp put/vacate cho đến khi thỏa.

\subsection{Đảm bảo không positive cycle}

Khi thêm max lag $(j, i, -L)$, kiểm tra positive cycle bằng max-plus closure
trước khi chấp nhận cạnh:

\begin{verbatim}
def add_max_lag_safe(lags, i, j, ell_min, slack):
    L   = ell_min * (1 + slack)
    candidate = (j, i, -L)
    if not has_positive_cycle(lags + [candidate]):
        lags.append(candidate)
    # else: bỏ qua max lag này
\end{verbatim}

Instance có positive cycle được giữ lại như một case UNSAT để kiểm tra
bộ phát hiện infeasibility của solver.

\subsection{Điều khiển độ khó}

\begin{table}[htbp]
\centering
\caption{Tham số điều khiển độ khó và ảnh hưởng lên solver.}
\label{tab:difficulty}
\begin{tabularx}{\textwidth}{@{}p{0.22\textwidth}p{0.09\textwidth}Y@{}}
\toprule
\textbf{Tham số} & \textbf{Tăng $\to$} & \textbf{Ảnh hưởng lên solver}\\
\midrule
$n_\mathrm{copies}$ & Khó hơn &
  Horizon tăng; nhiều symmetry cần breaking.\\
$P_s$ (storage cap) & Dễ hơn &
  Storage constraints lỏng hơn; ít conflicts hơn.\\
$\rho_\mathrm{max}$ (tỷ lệ max lag) & Dễ hơn &
  Windows hẹp hơn; formula nhỏ hơn.\\
$\delta_\mathrm{slack}$ (lag slack) & Khó hơn &
  Max lag xa min lag; windows rộng hơn; nhiều SAT calls hơn.\\
$\lambda_\mathrm{lag}$ (lag density) & Dễ hơn &
  Closure mạnh hơn; windows hẹp hơn; PB(AMO) lợi hơn.\\
$\sigma_\mathrm{storage}$ (storage density) & Khó hơn &
  Nhiều storage PB constraints hơn.\\
$B_r = 1$ (unary) & Cố định &
  No-overlap cấu trúc; path cover/PB(AMO) hiệu quả.\\
\bottomrule
\end{tabularx}
\end{table}

Ba cấu hình chuẩn:

\begin{table}[htbp]
\centering
\caption{Ba preset độ khó cho benchmark SAMSP-Synthetic.}
\label{tab:presets}
\begin{tabularx}{\textwidth}{@{}p{0.13\textwidth}ZZZZZZ@{}}
\toprule
\textbf{Preset} & $n_\mathrm{copies}$ & $n_\mathrm{types}$
  & $P_s$ & $\rho_\mathrm{max}$ & $\sigma_\mathrm{storage}$ & $\delta$\\
\midrule
Easy   & 1--5   & 1   & 4   & 0.10 & 0.30 & 0.5--2.0\\
Medium & 5--20  & 1--2 & 2  & 0.30 & 0.50 & 0.2--1.0\\
Hard   & 20--120 & 2--4 & 1 & 0.50 & 0.60 & 0.1--0.5\\
\bottomrule
\end{tabularx}
\end{table}

\subsection{Pipeline sinh instance}

\begin{enumerate}
  \item Sinh chuỗi backbone (chain) cho mỗi test type;
    thêm shortcut min lags theo $\lambda_\mathrm{lag}$.
  \item Thêm max lags an toàn (kiểm positive cycle từng cạnh).
  \item Assign renewable resources; assign storage events theo cặp.
  \item Ghép instance từ copies; thêm symmetry edges SB-R và SB-C.
  \item Max-plus closure toàn cục; tính $LB_0$; chạy greedy $\to$ $\overline{U}$.
  \item Tính windows; kiểm $\ES_i \le \LS_i$ cho mọi $i$.
  \item Kiểm storage feasibility trên greedy schedule; fix nếu vi phạm.
  \item Ghi metadata, seed và SHA-256 checksum.
\end{enumerate}

\subsection{Kế thừa từ \texttt{RCPSPProblem} hiện tại}

\begin{verbatim}
class SAMSPProblem(RCPSPProblem):
    """Mở rộng RCPSPProblem với generalized lags và storage."""

    def __init__(self, *args,
                 lags=None,            # list of (i, j, ell)
                 storages=None,        # list of storage names
                 storage_caps=None,    # list of P_s
                 storage_deltas=None,  # c[i][s]
                 test_metadata=None):
        super().__init__(*args)
        self._lags           = lags or []
        self._storages       = storages or []
        self._storage_caps   = storage_caps or []
        self._storage_deltas = storage_deltas or []
        self._test_metadata  = test_metadata or {}

    @classmethod
    def from_rcpsp(cls, rcpsp) -> 'SAMSPProblem':
        """Nâng cấp RCPSP thành SAMSP: không storage, không max lag.
        Dùng làm regression test trên PSPLIB."""
        lags = [(i, j, rcpsp.durations[i])
                for i, j in rcpsp.precedence_graph.edges()]
        return cls(
            number_of_activities=rcpsp.number_of_activities,
            number_of_resources=rcpsp.number_of_resources,
            durations=rcpsp.durations,
            requests=rcpsp.requests,
            capacities=rcpsp.capacities,
            lags=lags,
            storages=[],
            storage_caps=[],
            storage_deltas=[[0]] * rcpsp.number_of_activities,
        )
\end{verbatim}

\subsection{Phân tích instance trong bài báo gốc và hệ quả cho generator}
\label{sec:paper-instance}

\paragraph{Cấu trúc máy (cố định cho mọi instance).}
Ba renewable resources $\mathcal{R} = \{\text{Shuttle}, \text{Sample Arm}, \text{Reagent Arm}\}$, mỗi tài nguyên
capacity 1 ($B_r = 1$, unary).
Bốn storage areas $\mathcal{C} = \{\text{Cold}, \text{Hot-I}, \text{Hot-II}, \text{Observation}\}$,
capacity $P = \langle 4, 4, 4, 2 \rangle$.

\begin{remark}
Observation area có capacity $P_s = 2$ thay vì 4 — đây là ràng buộc thắt cổ chai.
Generator cần phản ánh sự bất đối xứng này thay vì dùng capacity đồng nhất.
\end{remark}


\paragraph{Bốn test types (Table~1 trong bài báo).}

\begin{table}[htbp]
\centering
\caption{Đặc trưng 4 test types. LB = cận dưới max-plus cho một copy đơn.}
\label{tab:testtypes}
\begin{tabularx}{\textwidth}{@{}p{0.09\textwidth}ZZZZZZZZZZ@{}}
\toprule
\textbf{Type} & \textbf{Acts} & \textbf{Prec.} & \textbf{LB} &
  \textbf{Shuttle} & \textbf{S.Arm} & \textbf{R.Arm} &
  \textbf{Cold} & \textbf{Hot-I} & \textbf{Hot-II} & \textbf{Obs.}\\
\midrule
1 &  7 &  8 & 25 & 3 & 1 & 1 & 0 & 1 & 0 & 1\\
2 & 14 & 17 & 62 & 6 & 5 & 1 & 2 & 0 & 1 & 1\\
3 &  9 & 11 & 44 & 4 & 1 & 2 & 0 & 1 & 1 & 1\\
4 &  9 & 10 & 38 & 4 & 1 & 2 & 1 & 0 & 1 & 1\\
\bottomrule
\end{tabularx}
\end{table}

Bốn test types phân biệt theo chuỗi storage zones mà cuvette đi qua:
Type~1 (Hot-I $\to$ Obs), Type~2 (Cold$\times$2 $\to$ Hot-II $\to$ Obs, hai cuvettes song song),
Type~3 (Hot-I $\to$ Hot-II $\to$ Obs), Type~4 (Cold $\to$ Hot-II $\to$ Obs).
Ba đặc điểm chung: (i)~mỗi activity chiếm đúng một renewable resource;
(ii)~time lag incubation lớn hơn duration thao tác khoảng 10 lần;
(iii)~transfer cuvette giữa hai storage zones được mô hình hóa trong một activity.

\paragraph{Quy mô và cấu trúc instance.}

Bài báo ký hiệu $I_{[\langle i, j \rangle]}$ với $i$ là test type, $j$ là số copies.
Set~I ($j \in \{5, 10, 20\}$) và Set~II ($j \in \{20, 60, 120\}$, kể cả mix nhiều types)
gồm 24 instances mỗi set.

\begin{table}[htbp]
\centering
\caption{Quy mô Set~I theo số copies (số liệu từ bài báo gốc).}
\label{tab:set1stats}
\begin{tabularx}{\textwidth}{@{}p{0.22\textwidth}ZZZY@{}}
\toprule
\textbf{Phạm vi} & \textbf{Acts} & \textbf{Prec. w/o SB} & \textbf{Prec. w SB} & \textbf{Nhận xét}\\
\midrule
Single type, 5 copies  & 27--62   & 131--541  & 341--1501   & Giải được dưới 1 giây.\\
Single type, 10 copies & 52--122  & 261--1081 & 1206--5401  & Thách thức trung bình.\\
Single type, 20 copies & 102--242 & 521--2161 & 4511--20401 & Khó nhất trong Set~I.\\
Mix 2 types, 5+5       & 62--97   & 336--791  & 866--2131   & Phức tạp hơn do asymmetry.\\
Mix 2 types, 10+10     & 122--192 & 671--1581 & 3146--7836  & Một số solver timeout.\\
\bottomrule
\end{tabularx}
\end{table}

Symmetry breaking làm tăng số precedences 3--10 lần và đồng thời cải thiện $LB$
(ví dụ: Type~1, 5 copies --- prec tăng từ 131 lên 341, $LB$ tăng từ 25 lên 59).
SB edges phải được tích hợp vào closure \emph{trước} khi tính time windows.
Set~II có các instance chưa giải được với mọi solver;
SAMSP-Synthetic-Hard nhắm đến 100--400 activities để tái tạo vùng khó tương ứng.

\paragraph{Cài đặt generator.}

Generator tái tạo đúng cấu trúc 4 test types qua 5 điều chỉnh so với phiên bản ban đầu:
(1)~hỗ trợ transfer cuvette giữa hai storage zones trong một activity;
(2)~capacity không đồng nhất giữa các storage areas (Observation = 2 < Cold/Hot = 4);
(3)~tham số lag incubation độc lập với duration thao tác;
(4)~sửa lỗi indexing khiến instance ghép nhiều test types bị tính thêm dummies thừa và trả LB sai;
(5)~số activities khớp chính xác Table~1 (5/12/7/7 real activities per type).

\begin{table}[htbp]
\centering
\caption{So sánh generator với bài báo (Table~2 gốc).}
\label{tab:generator-vs-paper}
\begin{tabularx}{\textwidth}{@{}lZZZZZZ@{}}
\toprule
\textbf{Instance} & \multicolumn{2}{c}{\textbf{Acts}} &
  \multicolumn{2}{c}{\textbf{Prec w/o SB}} &
  \multicolumn{2}{c@{}}{\textbf{LB w SB}}\\
 & Ours & Paper & Ours & Paper & Ours & Paper\\
\midrule
$I_{\{1,5\}}$  & \textbf{27}  & 27  &  71 &  131 & 49 & 59\\
$I_{\{1,10\}}$ & \textbf{52}  & 52  & 226 &  261 & 80 & 100\\
$I_{\{1,20\}}$ & \textbf{102} & 102 & 800 &  521 & 150 & 185\\
$I_{\{2,5\}}$  & \textbf{62}  & 62  & 115 &  541 & 82 & 128\\
$I_{\{3,5\}}$  & \textbf{37}  & 37  & 121 &  251 & 60 &  88\\
$I_{\{4,5\}}$  & \textbf{37}  & 37  &  80 &  206 & 58 &  66\\
\bottomrule
\end{tabularx}
\end{table}

Số activities khớp 100\% với bài báo. LB w~SB thấp hơn do symmetry breaking theo
capacity storage chưa tái tạo đúng thứ tự cuvette theo thời gian --- đây là hướng cải thiện tiếp theo.

\paragraph{Trạng thái benchmark.}

\begin{table}[htbp]
\centering
\caption{Benchmark sets. Repo: \texttt{github.com/kieuvantuyen01/samsp-sat-solver}.}
\label{tab:benchmarks-updated}
\begin{tabularx}{\textwidth}{@{}p{0.26\textwidth}ZYY@{}}
\toprule
\textbf{Set} & \textbf{N} & \textbf{Acts} & \textbf{Trạng thái}\\
\midrule
SAMSP Set~I      & 24  & 27--242   & \checkmark\ \texttt{data/set\_i/}\\
SAMSP Set~II     & 24  & 97--1442  & \checkmark\ \texttt{data/set\_ii/}\\
Synthetic-Easy   & 100 & 10--50    & \checkmark\ \texttt{--preset easy}\\
Synthetic-Medium & 100 & 50--300   & \checkmark\ \texttt{--preset medium}\\
Synthetic-Hard   &  50 & 200--1600 & \checkmark\ \texttt{--preset hard}\\
SAMSP-UNSAT      &  20 & 6+        & \checkmark\ \texttt{--unsat}\\
RCPSP-Regression & 960 & 32--62    & planned\\
\bottomrule
\end{tabularx}
\end{table}


\appendix
\section{Bảng ký hiệu}
%=============================================================================

\begin{longtable}{@{}p{0.22\textwidth}p{0.70\textwidth}@{}}
\toprule \textbf{Ký hiệu} & \textbf{Ý nghĩa}\\ \midrule
\endfirsthead \toprule \textbf{Ký hiệu} & \textbf{Ý nghĩa}\\ \midrule \endhead
$\mathcal{V},\, \mathcal{A}$ & Tập mọi hoạt động; tập hoạt động không dummy.\\
$d_i$ & Thời lượng.\\
$\ell_{ij},\, \ell^\star_{ij}$ & Lag gốc; lag mạnh nhất từ max-plus closure.\\
$\ES_i, \LS_i, \EC_i, \LC_i$ & Earliest/latest start/completion.\\
$W_i = \LS_i - \ES_i + 1$ & Độ rộng start window.\\
$U_i = W_i + d_i - 1$ & Độ rộng run window.\\
$q_{i,t}$ & Order literal: $A_i$ đã bắt đầu không muộn hơn $t$.\\
$\sigma_{i,t} = q_{i,t}\wedge\neg q_{i,t-1}$ & Start literal dẫn xuất.\\
$u_{i,t}$ & Run literal: $A_i$ đang chạy tại $t$.\\
$b_{ir},\, B_r$ & Nhu cầu và sức chứa tài nguyên tái tạo.\\
$c_{is},\, P_s$ & Thay đổi level và sức chứa storage.\\
$O_s(t)$ & Mức chiếm storage $s$ tại biên $t$.\\
$D = (\mathcal{A}, F)$ & DAG non-overlap dùng cho PB(AMO).\\
$I_s^+,\, I_s^-$ & Activities chiếm / trả storage $s$.\\
\bottomrule
\end{longtable}

%=============================================================================
\section{Checklist trước khi báo cáo}
%=============================================================================

\begin{enumerate}
  \item Mọi model SAT giải mã được đều qua validator kiểm \eqref{eq:dummy}--\eqref{eq:storage}.
  \item Optimum khớp baseline trên instance nhỏ và benchmark đã biết.
  \item Có test riêng cho negative lag, positive-cycle UNSAT, zero-duration dummy.
  \item Ghi rõ PB encoder cụ thể; không gọi BDD PB thông thường là PB(AMO).
  \item Tách $\#V_\mathrm{core},\#C_\mathrm{core}$ khỏi $\#V_\mathrm{PB},\#C_\mathrm{PB}$.
  \item Cùng solver version, seed, timeout và phần cứng cho mọi cấu hình.
  \item Cung cấp script tái lập kết quả và kiểm tra schedule.
\end{enumerate}

%=============================================================================
\begin{thebibliography}{9}
%=============================================================================

\bibitem{bofill2020rcpsp}
M.~Bofill, J.~Coll, J.~Suy, and M.~Villaret.
SMT encodings for resource-constrained project scheduling problems.
\emph{Comput.\ Ind.\ Eng.}, 149:106777, 2020.

\bibitem{bofill2022samsp}
M.~Bofill, J.~Coll, G.~Martín, J.~Suy, and M.~Villaret.
The Sample Analysis Machine Scheduling Problem.
\emph{Comput.\ Oper.\ Res.}, 142:105730, 2022.

\bibitem{bofill2021pbamo}
M.~Bofill et al.
SAT encodings for pseudo-Boolean constraints together with at-most-one constraints.
\emph{Artif.\ Intell.}, 302:103604, 2021.

\bibitem{truong2025scl}
H.~X.~Truong, T.~V.~Kieu, and K.~V.~To.
Sequential counter encoding for staircase at-most-one constraints.
Manuscript, 2025.

\bibitem{abio2012bdd}
I.~Abío et al.
A new look at BDDs for pseudo-Boolean constraints.
\emph{J.\ Artif.\ Intell.\ Res.}, 45:443--480, 2012.

\bibitem{kolisch1997psplib}
R.~Kolisch and A.~Sprecher.
PSPLIB -- A project scheduling problem library.
\emph{Eur.\ J.\ Oper.\ Res.}, 96(1):205--216, 1997.

\end{thebibliography}

\end{document}
